\bookchapter{附录A：新手机开箱“保姆级”流程}

\booksection{开箱验机“三部曲”：}

\begin{itemize}
  \item \textbf{看外观}：检查机身、屏幕有无划痕、磕碰。
  \item \textbf{对信息}：核对盒子上的IMEI/序列号和发票上的是否一致。
  \item \textbf{试功能}：拨个电话试试听筒话筒，拍张照片检查摄像头。
\end{itemize}

\booksection{激活与初始设置“重点提醒”：}

\begin{itemize}
  \item \textbf{网络}：先连上稳定的Wi-Fi。
  \item \textbf{账号}：登录你的苹果ID、华为账号、小米账号等。
  \item \textbf{基础设置}：把"查找我的手机"打开，设置一个靠谱的锁屏密码。
  \item \textbf{eSIM设置}：如果你的手机支持eSIM（如iPhone、部分安卓旗舰），可以在设置中添加eSIM套餐，无需插入实体SIM卡。出国时eSIM尤其方便，可以直接购买当地运营商的eSIM套餐。
  \item \textbf{数据迁移}：从旧手机迁移数据时，确保两台手机电量充足。iPhone用"快速开始"，安卓用官方换机App。
\end{itemize}

\booksection{装App“安全第一”：}

\begin{itemize}
  \item \textbf{首选官方商店}：苹果用App Store，安卓用手机自带的应用商店。
  \item \textbf{权限别乱给}：安装时弹出的权限请求，想清楚"这个App真的需要吗？"
\end{itemize}
